A \code{ThreadPool} is a collection of threads that make up a
process. Each \code{Process} has one \code{ThreadPool}. A \code{ThreadPool}
is typically used to iterate over or search the set of Threads. It satisfies the C++
\href{https://en.cppreference.com/cpp/named_req/ForwardIterator}{LegacyForwardIterator}
concept. It also offers the standard \code{begin/end} members for iteration.

Note that it is not safe to make assumptions about having consistent contents of
a ThreadPool for a running target process. As the target process runs, threads
may be inserted or removed from the ThreadPool. It is generally safer
to stop a Process before operating on its ThreadPool. When used on a running
process the ThreadPool iterator methods guarantee that they will not return
invalid Thread objects (e.g, nothing that would lead to a segfault), but they
do not guarantee that the Thread objects will refer to live threads or that
they return all Threads.

\begin{apient}
ThreadPool::iterator find(Dyninst::LWP lwp)
ThreadPool::const_iterator find(Dyninst::LWP lwp) const
\end{apient}

\apidesc{
The functions respectively return an iterator and a const\_iterator that points
to the Thread with a LWP of lwp. If lwp is not found in the thread list, then
this function returns end().
}

\begin{apient}
size_t size() const
\end{apient}

\apidesc{
This function returns the number of Threads in the Process.
}

\begin{apient}
Process::ptr getProcess()
Process::const_ptr getProcess() const 
\end{apient}

\apidesc{
These functions respectively return a pointer or a const pointer to the Process
that owns this ThreadPool.
}

\begin{apient}
Thread::ptr getInitialThread()
Thread::const_ptr getInitialThread() const
\end{apient}

\apidesc{
These functions respectively return a pointer or a const pointer to the initial
Thread in a Process. The initial thread is the thread that started execution
of the process (i.e., the thread that called main).
}